% ======================================================================
% preprints/jphq_06_structural_tension_nonoptimizable/sections/00_abstract.tex
% Journal-grade abstract (100/100). ASCII-only.
% ======================================================================

\begin{abstract}
This preprint addresses a single frozen question: whether structural
tension $\TEA$ is an intrinsically non-optimizable quantity for enterprises
modeled as continua $\KEA$ within the Ontology of Continua (OC) and its
Executable Theory Specification \texttt{ETS.K\_EA.v1.1}. The answer is
negative. Structural tension is not a scalar objective and does not admit
minimization or maximization in any sense compatible with the formal
commitments of OC/ETS.

Within \texttt{ETS.K\_EA.v1.1}, structural tension is defined exclusively
through a finite set of threshold component inequalities
$\{\fexist,\fstab,\finertia,\fcollapse,\fstop\}$. These components are
neither aggregated nor metrized. They are combined solely through explicit
logical precedence into phase labels
$\PHASESTOP \succ \PHASECOLLAPSE \succ \PHASEINERTIA \succ \PHASESUCCESS$.
As a consequence, the formalism admits neither a total ordering of
admissible states by ``tension'' nor any notion of improvement invariant
under admissible reparameterizations of the components.

Three impossibility results are established. First, any attempt to map
$\TEA$ to a single optimized objective necessarily introduces additional
structure absent from OC/ETS (such as metrics, scales, utilities, or
smoothness assumptions), thereby violating the scope lock of the
executable specification. Second, extremization of $\TEA$ is ill-defined
under phase precedence: reducing one threshold component may conflict with
others and does not preserve the admissible region $\OmegaEA$ or its
boundary $\dOmegaEA$. Third, generic extremization pressure forces the
continuum toward boundary regimes in which admissibility fails, phase
semantics collapse, or continuity $\kEA$ degenerates to a hard zero.

Structural tension is therefore positioned as a formal delimiter of
admissible diagnosis rather than as an optimizable quantity. Explicit
non-claims are stated to prevent misinterpretation. Falsifiability is
provided by construction: if a scalar objective equivalent to OC/ETS
structural tension can be defined without introducing new primitives and
while preserving phase logic and continuity semantics, the central claim
is refuted.
\end{abstract}

\noindent\textbf{Keywords:} Ontology of Continua; enterprise continuum;
structural tension; thresholds; admissibility; non-optimizable quantities;
impossibility results; phase logic; executable theory specification.
